\documentclass[11pt]{jsarticle}

\title{ターン制戦略ゲームにおける\\味方NPCの戦術生成による振る舞いの違い}
\author{2023TS023：可児桜将}
\date{2025年12月21日}

\begin{document}

\maketitle

\section{研究背景}
近年、ゲームAIは大きな発展を遂げており、敵キャラクタの強化だけでなく、プレイヤーの行動やプレイスタイルを考慮したNPCについての研究も進められている。しかしその一方で、プレイヤーの味方として行動するNPCに対しては「役に立たない」「戦術的に不自然である」といった不満の声が多く見られる。

特にターン制戦略ゲームでは、味方NPCの行動がプレイヤーの戦略と噛み合わず、ゲーム体験を損なう要因となることが指摘されている。しかし、敵AIと比較して味方NPCに焦点を当てた研究は少なく、十分な検証が行われていないのが現状である。
\subsection{研究目的}
本研究の目的は、ターン制戦術ゲームにおける味方NPCの行動に着目し、戦術生成アルゴリズムの違いがNPCの振る舞いやゲーム結果にどのような影響を与えるかを明らかにすることである。

\section{ゲーム環境}
実験のため、市販のターン制戦略ゲーム及び既存研究を参考に独自のゲーム環境を構築した。ステージは正方形のマスで構成されたマップであり、マスごとに移動コストや通行不可の設定が存在する。

コマは攻撃範囲や移動可能マスの異なる4種類を用意し、自軍（プレイヤー）、敵軍（NPC）、友軍（味方NPC）の3勢力が存在する。

\subsection{実験方法}
実験には、市販ゲームにおいて不マンの声が多く上がっていたステージをもとに作成したマップを使用した。NPCの行動には、既存研究及び解析によって知られている戦術アルゴリズムを用いた

実験条件は以下の2通りとし、それぞれ100回の対戦を行った。
\begin{itemize}
\item 敵軍のみに戦術アルゴリズムを使用した場合
\item 敵軍及び友軍の療法に戦術あるごリズっむを適用した場合
\end{itemize}

これらの結果を比較することで、味方NPCへのアルゴリズム適用がゲーム結果に与える影響を分析した。

\section{実験結果}
実験の結果、敵軍のみにアルゴリズムを適用した場合には自軍の勝利数が多く、敵軍の勝利はほとんど見られなかった。一方で、敵軍と友軍の両方にアルゴリズムを適用した場合には、敵軍の勝利が増加し、自軍の勝率が大きく低下した。

この結果から、味方NPCに戦術アルゴリズムを適用することが、必ずしもプレイヤーに撮った有利に働かないことが明らかとなった。

\section{考察}
同一の戦術生成アルゴリズムであっても、敵NPCと味方NPCでは求められる役割が異なる。敵NPCは合理的かつ最適な行動を取ることでプレイヤーに挑戦を与える存在であるのに対し、味方NPCはプレイヤーの戦略を補助する存在である。

しかし本研究の結果から、味方NPCが合理的な行動を優先することによって、かえってプレイヤーの糸を妨害し、結果的に敵を有利にしてしまう可能性が示唆された。これは、味方NPCに求められる要素が単なる強さではなく、協調性や支援性であることを示している。

\section{まとめと今後の課題}
本研究では、ターン制戦略ゲームにおいて既存の戦術生成アルゴリズムを味方NPCに適用した場合、ランダム行動よりも自軍の勝率が低下する場合があることを示した。

今後は、戦術アルゴリズムとランダム行動を組み合わせることで、プレイヤーの意図を尊重しつつ協力的に行動する味方NPCの実現を目指す必要がある。また、プレイヤーの行動や戦況を考慮した柔軟な戦術生成手法の検討が今後の課題である

\begin{thebibliography}{99}
\bibitem{ref1}
  野戸彰大, 今井俊之,
  ターン制戦略ゲームの味方NPCの戦術生成部分による振舞いの違い,
  情報処理学会第86回全国大会, 2024.
\end{thebibliography}

\end{document}
